\pdfoutput=1
\ifdefined\pdfinfoomitdate\pdfinfoomitdate=1\fi
\ifdefined\pdftrailerid\pdftrailerid{}\fi
\documentclass{ledger}
\usepackage{ledger-manuscript}
\usepackage{booktabs,longtable,array,mathtools}
\usepackage{tikz}
\usetikzlibrary{arrows.meta,positioning,calc,fit}
\usepackage{flafter,placeins}
\hypersetup{
  pdftitle={GL1F: Reproducible Integer Tree-Ensemble Inference on the EVM},
  pdfauthor={Mikhail Fedorov},
  pdfsubject={Author manuscript dated 6 September 2026; not peer reviewed},
  pdfkeywords={tree ensembles, reproducibility, EVM, integer inference}
}
\setlength{\emergencystretch}{2em}
\newcommand{\Ithirtytwo}{\mathbb{I}_{32}}
\newcommand{\codeword}[1]{\texttt{#1}}
\newcommand{\hashpair}[2]{\texttt{#1}\allowbreak\texttt{#2}}
\title{GL1F: Reproducible Integer Tree-Ensemble Inference on the EVM}
\author{Mikhail Fedorov\thanks{Corresponding author.}\thanks{M. Fedorov is the founder of GenesisL1.}}

\begin{document}
\maketitle
\thispagestyle{pagefirst}

\begin{abstract}
Public model storage does not by itself establish which function a registered
model executes. We study this gap for fixed-depth integer tree ensembles on
the Ethereum Virtual Machine (EVM). GL1F combines a canonical scalar/vector
encoding, ordered code-as-data chunks, and separate local and contract
evaluators. We specify sufficient conditions for exact reconstruction and
integer execution agreement, including registry/header consistency and
boundary-spanning reads. An executable counterexample shows that registration
alone does not bind an identifier to the referenced bytes. JavaScript, Python,
and C++ trainers emit identical cores for 25 tested profiles; arithmetic
counterexamples delimit that empirical result. Across six scalar shapes, code-as-data publication uses 61.8--69.4\% less
gas than packed Solidity storage, with mixed inference costs. A pinned-state
archive makes 108 output comparisons for 12 models totaling 31,185,324 bytes
replayable without a live provider. The resulting contribution is a reproducible method
for checking the connection between a public model identifier, its bytes, and
its executed integer function. It establishes conditional execution agreement,
without establishing predictive validity or authenticating provider-reported
chain history.
\end{abstract}

\section{Introduction}
\label{sec:introduction}

Publishing a trained model does not fully specify a reproducible prediction.
A result also depends on feature ordering, preprocessing, numeric conventions,
serialization, and the implementation that interprets the artifact.
Reproducibility initiatives in computational science and machine learning
therefore emphasize executable artifacts and explicit evaluation protocols.
\cite{peng2011,pineau2021} A blockchain supplies a persistent, ordered record
of public data and state transitions, but the association between a model
identifier and an executed function still requires verification. A registry
can contain inconsistent metadata; a decoder can interpret malformed bytes
differently; and a remote procedure call (RPC) can report a result without
providing evidence that authenticates the underlying state.

This paper addresses a bounded version of that problem: under which
representation, storage, state, input, and arithmetic conditions do local and
EVM evaluators produce exactly the same integer result for a registered tree
ensemble? The question is relevant when a contract must consume the model
output directly, or when researchers need to reconstruct the precise public
function available at an earlier block. For applications requiring only local
arithmetic replay, a content-addressed file and a specified evaluator can
suffice. The additional role of the ledger is to order public artifacts and
service state and make them accessible to contract execution.

Gradient-boosted decision trees provide a useful restricted model family.
They are widely used for tabular learning,\cite{friedman2001,chen2016xgboost,ke2017lightgbm}
but their inference kernel needs only comparisons, indexed leaf selection,
and addition. Complete fixed-depth trees make both offsets and work counts
explicit. Integer thresholds, leaves, and packed inputs avoid floating-point
ambiguity during inference. These properties support a small execution
specification, while retaining scalar scores and vector logits for several
training objectives.

GL1F implements that design with browser, Python, and C++ trainers, a
canonical binary model format, and a Solidity interpreter. Models exceeding a
single contract's code capacity are split across public code-as-data objects.
An ordered address table and registry record identify the storage layout.
The implementation is deployed on GenesisL1, an EVM-compatible network. The
model bytes are intentionally public, allowing reconstruction and local
inference even when a canonical service endpoint is metered.

The contribution is the composition of three checkable relations. First, an
ordered storage manifest must reconstruct one exact canonical byte string.
Second, the registry fields actually consumed by the runtime must agree with
that string. Third, evaluators must apply the same signed decoding,
strict-greater branch rule, and exact accumulation to the same packed input.
We state sufficient conditions for these relations and provide a short
specification-level equivalence proof. An executable counterexample
establishes why the registered identifier additionally needs an explicit
content-hash check. Cross-language and local-EVM tests then assess the
implemented paths, including numeric and chunk-boundary failures. A
controlled storage experiment and a pinned deployment archive evaluate the
cost and practical observability of this construction.

Exact on-chain tree inference, quantized parameters, and code-as-data storage
are established techniques.\cite{alhaidari2025onchain,sequence_sstore2}
The present study examines the connection from a registry-named collection of
objects to its executed integer function, rather than claiming priority for
those components. Its results also distinguish two forms of reproducibility:
inference equality follows under stated premises, whereas training equality
is supported only for the executed profiles and numeric environment. The
arithmetic counterexamples are informative precisely because they prevent a
finite test matrix from being mistaken for a universal training guarantee.

The evaluation concerns execution assurance. It does not test whether the
published models predict their application targets accurately. Throughout,
``independent evaluator'' means a separately implemented checking path;
the author developed the system and conducted the experiments, so this is not
an externally replicated study. This distinction permits specific, falsifiable
claims about representation and execution while leaving model quality and
chain-state authentication to their own evidence.

\section{Related work}
\label{sec:related}

\subsection{Portable and blockchain-executed models}

Gradient boosting constructs additive predictors by fitting successive weak
learners to a loss-derived target.\cite{friedman2001} XGBoost and LightGBM
provide mature training systems with regularization, histogram methods, and
scalable execution.\cite{chen2016xgboost,ke2017lightgbm} GL1F restricts tree
shape for explicit offsets and bounded inference work; the experiments here
do not compare its predictive performance with those systems.

ONNX already defines a portable tree-ensemble operator, including branch
modes, leaf targets, aggregation, and post-transforms.\cite{onnx_treeensemble}
Hummingbird compiles conventional machine-learning pipelines to tensor
computations for prediction serving.\cite{nakandala2020hummingbird} GL1F's
smaller grammar provides a fixed-width integer object that can be interpreted
directly by the EVM. The tradeoff is reduced model and operator flexibility,
not a general replacement for model interchange.

Prior blockchain systems have evaluated collaborative model updating on
Ethereum, logistic-regression inference in Fabric, integer Naive Bayes in
Solidity, and off-chain-trained bagged trees deployed as prediction
contracts.\cite{harris2019decentralized,drungilas2021federated,badruddoja2022predict,kadadha2022behavior}
ML2SC translates PyTorch multilayer perceptrons to Solidity, transfers their
parameters, applies fixed-point arithmetic, and studies gas costs.
\cite{li2024ml2sc} These results establish that off-chain training followed by
contract inference is not specific to GL1F.

The closest property-level overlap is Alhaidari et al., who report quantized,
serialized, hash-bound model parameters, read-call and transaction inference,
gas analysis, and reported Z3 verification of bit-exact model execution,
with decision-tree support.\cite{alhaidari2025onchain} Consequently, bit-exact tree
execution and parameter hash binding cannot serve as standalone novelty
claims here. GL1F instead supplies a strict reconstruction predicate for an
ordered multi-object artifact, reconciles it with the metadata consumed by a
runtime, and demonstrates a concrete case in which a registry key fails to
bind the referenced model. Its scalar/vector grammar, conformance matrix, and
replayable deployment record make that composition inspectable.

\subsection{Storage, verification, and provenance}

The EVM exposes deterministic execution and resource charging through gas.
\cite{wood2014ethereum} EIP-170 limits deployed runtime code to 24,576 bytes.
\cite{buterin2016eip170} Code-as-data libraries such as SSTORE2 already use
contract creation to publish bytes and \codeword{EXTCODECOPY} to recover
them.\cite{sequence_sstore2} The storage primitive is therefore prior art.
Here the relevant object is the ordered table of chunk addresses, exact
framing and lengths, and the boundary reader used by a model interpreter.
The controlled experiment isolates a practical tradeoff between that design
and a packed Solidity-storage implementation.

Fu et al. authenticate binary decision-tree predictions using committed
trees and provider-supplied path proofs.\cite{fu2023vdt} Zero-knowledge
systems prove decision-tree predictions or accuracy while protecting selected
private information,\cite{zhang2020decisiontreezk,singh2022zkpipelines}
whereas opML uses off-chain execution and an optimistic dispute protocol.
\cite{conway2024opml} GL1F directly interprets public weights. It requires no
per-prediction prover or challenge window, but it incurs replicated execution
and provides no model confidentiality. The experiments do not rank those
architectures on a common workload.

KEVM gives a formal semantics of EVM bytecode.\cite{hildenbrandt2018kevm}
The theorem in this paper relates conforming evaluator specifications; it is
not a KEVM proof, a verified compiler translation, or a mechanized proof of
the deployed bytecode. Source inspection and finite execution tests support
the separate implementation-conformance claim.

Supply-chain provenance systems address additional questions. in-toto
links authorized processing steps and their materials,\cite{torresarias2019intoto}
while MLflow2PROV records provenance across machine-learning pipelines.
\cite{schlegel2025provenance} A content digest in GL1F identifies bytes but
does not identify their trainer, establish dataset lineage, or authenticate a
scientific interpretation. Likewise, a fixed compiler configuration does not
by itself demonstrate reproducible builds.\cite{lamb2022reproducible}

The public GL1F protocol document dated 19 January 2026 already describes the
alpha format, storage architecture, contracts, and browser workflow.
\cite{gl1fprotocol2026} The present article evaluates that system with
explicit verification conditions, a registry counterexample, adversarial
cross-language conformance tests, and quantitative storage and deployment
evidence. The earlier document remains identifiable as prior system
documentation.

\section{Representation and system design}
\label{sec:design}

\subsection{Canonical integer forest}

Let $\Ithirtytwo=[-2^{31},2^{31}-1]\cap\mathbb Z$. A complete binary tree of
depth $d$ has $I(d)=2^d-1$ internal nodes and $L(d)=2^d$ leaves. Each internal
node occupies eight bytes: an unsigned 16-bit feature index, a signed 32-bit
threshold, and two zero reserved bytes. Each leaf is a signed 32-bit
increment. Multibyte fields are little-endian, and nodes use flat heap order.
The tree length is therefore
\begin{equation}
 b(d)=8(2^d-1)+4\cdot2^d=12\cdot2^d-8.
 \label{eq:treebytes}
\end{equation}
A 24-byte fixed header records the version, feature count, depth, tree count,
base or output count, and positive quantization scale $Q$.

Version 1 represents a scalar ensemble with base $a_0\in\Ithirtytwo$ and $T$
trees. Version 2 represents $K\ge2$ outputs, $R\ge1$ trees per output, and a
base vector in $\Ithirtytwo^K$. Its trees are output-major, and its base
vector follows the fixed header. The canonical lengths are
\begin{equation}
 |B_{\mathrm{v1}}|=24+Tb(d),\qquad
 |B_{\mathrm{v2}}|=24+4K+KRb(d).
 \label{eq:lengths}
\end{equation}
These are representation identities rather than statistical model-size
estimates. A canonical parser checks the magic, recognized version, zero
reserved fields, positive feature count and scale, exact length, supported
depth, and every feature index. The evaluated deployment profile uses
$1\le d\le12$. For a record satisfying registry/core agreement, the registry's 16-bit
total-tree field requires $T\le65535$ for v1 and $KR\le65535$ for v2. Those bounds are narrower than the
32-bit tree-count fields in the core. Version 1 admits a base-only object
with $T=0$; trained evaluation fixtures use positive tree counts. Exact
header offsets appear in the supplement and the repository format
specification.

An input is exactly $4F$ bytes containing $q\in\Ithirtytwo^F$. It has already
been quantized and ordered before it reaches the inference interface. Starting
at heap index $i_0=0$, each decision is
\begin{equation}
 i_{j+1}=2i_j+1+\mathbf 1[q_{f(i_j)}>\tau(i_j)].
 \label{eq:traversal}
\end{equation}
Equality goes left. After $d$ decisions, the selected leaf is indexed by
$i_d-I(d)$. The scalar function and each vector component are
\begin{equation}
 f_B(q)=a_0+\sum_{t=0}^{T-1}v_{t,\pi_t(q)},\qquad
 f_{B,k}(q)=a_{0,k}+\sum_{t=0}^{R-1}v_{k,t,\pi_{k,t}(q)}.
 \label{eq:prediction}
\end{equation}
A class endpoint returns the lowest index among equal maximal integer logits.
Probabilities, calibration, and feature transformations belong to application
code, outside this packed-input function.

The core is an artifact identity, not a semantic normal form: distinct cores
can implement the same function. Conversely, identical scalar bytes may be
labelled as a regression score or a binary logit. The core does not encode
feature names, units, dataset provenance, or a training-task discriminator.
Optional \codeword{GL1X} JSON metadata is excluded from core hashing. A
scientific interpretation therefore needs a separately authenticated record of
preprocessing, feature and label order, and intended output meaning.

\subsection{Publication and reconstruction}

\codeword{ModelStore} publishes runtime code of the form
$\codeword{GL1C}\Vert\mathrm{payload}$. Its four-byte prefix leaves at most
24,572 payload bytes within EIP-170. A second code object contains an ordered
table of 32-byte address slots; the low 20 bytes identify a chunk and the high
12 bytes must be zero. One such table can hold at most
$\lfloor24572/32\rfloor=767$ addresses. For a core of length $\ell$, a nominal
payload stride $c$, and $N=\lceil\ell/c\rceil$ chunks, the final chunk may be
shorter than $c$; every earlier chunk must have exactly $c$ payload bytes.

The registry records the table pointer, byte length, stride, chunk count, and
inference fields. \codeword{ForestRuntime} reads signed and unsigned
primitives through \codeword{EXTCODECOPY}. A read crossing a chunk boundary
copies its prefix from one payload and its suffix from the next. Requiring
$c\ge4$ ensures that each one-, two-, or four-byte primitive spans at most
two chunks. The prefix \codeword{GL1C} is a framing marker, not an
authenticator: correctness also requires the exact table, order, lengths,
and payloads.

The publication workflow serializes a core, validates it locally, computes
Ethereum Keccak-256, stores its chunks and address table, and registers the
declared metadata. Verification reads those objects back at a selected state,
reconstructs the exact core, repeats canonical parsing, reconciles registry
fields, and checks the digest against the requested identifier. SHA-256 is
also recorded for ordinary file comparisons. Ethereum Keccak-256 here means
the legacy Ethereum hash, not FIPS SHA3-256.

\begin{figure}[t]
\centering
\resizebox{\linewidth}{!}{%
\begin{tikzpicture}[
  font=\small,
  node distance=8mm and 8mm,
  box/.style={
    draw=GLBlue!75!black,
    rounded corners=2pt,
    line width=0.7pt,
    fill=GLBlue!5,
    align=center,
    minimum height=10mm,
    text width=27mm
  },
  live/.style={
    box,
    draw=GLGreen!75!black,
    fill=GLGreen!6
  },
  verify/.style={
    box,
    draw=GLGold!80!black,
    fill=GLGold!8
  },
  flow/.style={-{Latex[length=2.4mm]}, line width=0.7pt, draw=GLInk!75},
  check/.style={-{Latex[length=2.4mm]}, line width=0.7pt, dashed, draw=GLGold!80!black}
]

\node[box] (browser) {Browser trainer\\and studio};
\node[box, below=of browser] (python) {Python CLI\\trainer};
\node[box, below=of python] (cpp) {C++17 CLI\\trainer};

\node[box, right=14mm of python, text width=31mm] (core) {Canonical GL1F core\\\texttt{GL1F} header +\\fixed-depth trees};
\node[verify, above=of core, text width=31mm] (oracle) {Strict independent parser\\hash and conformance report};
\node[verify, below=of core, text width=31mm] (local) {Local integer evaluator\\scalar or vector output};

\node[live, right=14mm of core] (chunks) {Immutable code-as-data\\\texttt{GL1C} chunks};
\node[live, above=of chunks] (table) {Ordered pointer table\\\texttt{GL1C} + addresses};
\node[live, below=of chunks] (registry) {Registry + Model NFT\\mutable service metadata};

\node[live, right=14mm of chunks, text width=29mm] (runtime) {ForestRuntime\\direct EVM integer\\inference};
\node[verify, above=of runtime, text width=29mm] (witness) {Pinned-state witness\\reconstruct, hash,\\compare outputs};
\node[box, below=of runtime, text width=29mm] (dapp) {Browser dApp or\\contract caller};

\draw[flow] (browser) -- (core);
\draw[flow] (python) -- (core);
\draw[flow] (cpp) -- (core);
\draw[flow] (core) -- (chunks);
\draw[flow] (chunks) -- (table);
\draw[flow] (table) -- (registry);
\draw[flow] (registry) -- (runtime);
\draw[flow] (chunks) -- (runtime);
\draw[flow] (dapp) -- (runtime);

\draw[check] (oracle) -- (core);
\draw[check] (core) -- (local);
\draw[check] (witness) -- (table);
\draw[check] (witness) -- (chunks);
\draw[check] (witness) -- (runtime);
\path (oracle.north) ++(0,7mm) coordinate (verifytop);
\draw[check] (local.west) -- ++(-6mm,0) |- (verifytop) -| (witness.north);

\node[
  draw=GLInk!35,
  rounded corners=2pt,
  fill=white,
  align=left,
  font=\footnotesize,
  below=13mm of registry,
  text width=91mm
] (boundary) {
\textbf{Assurance boundary.}
Green boxes are components observed at the pinned GenesisL1 deployment.
Gold boxes provide separate checking paths. Registry existence alone does not
authenticate a model: the witness must establish exact reconstruction,
canonical decoding, registry/header agreement, and
$\operatorname{keccak256}(\mathrm{core})=\mathrm{modelId}$.
};

\end{tikzpicture}
}
\caption{GL1F publication and inference path. Solid arrows show production
data flow; dashed arrows show separate verification paths. Model bytes are
intentionally public. Registry, NFT, pricing, and marketplace state are
services around the model and do not alter its mathematical inference
semantics.}
\label{fig:architecture}
\end{figure}


\subsection{Mutable service state}

The model NFT and optional marketplace surround the public artifact with
ownership, discovery, authorization, and payment services. These services
can change while the model bytes remain unchanged. In particular, token
ownership, endpoint access, and the fee destination are distinct registry
relations. An NFT transfer does not establish authorship or validate the
model's dataset. The evaluated NFT implements a custom ownership interface;
full ERC-721 metadata and discovery interoperability is not assumed.
\cite{eip721}

A paid endpoint can meter its own execution or service integration. It cannot
prevent anyone who reconstructs the public bytes from evaluating them
locally. A transaction exposes its packed input as public calldata. A view
call avoids that transaction but still reveals the request to the RPC
provider. Authorization therefore controls an endpoint rather than providing
input or model confidentiality.

\section{Execution assurance}
\label{sec:assurance}

\subsection{Sufficient conditions}

Consider one EVM state $\sigma$, registry record $r$, core $B$, and manifest
$M=(P,c,N,\ell)$. Call the manifest well formed when $P\ne0$,
$4\le c\le24572$, $1\le N\le767$, $\ell=|B|$, and
$N=\lceil\ell/c\rceil$; the table code is exactly $4+32N$ bytes with the
required prefix and canonical nonzero address slots; and chunk $j$ is exactly
$\codeword{GL1C}\Vert B[jc:\min((j+1)c,\ell)]$. All objects and registry
getters must refer to the same state. Exact final lengths exclude both
truncation and undeclared trailing payload.

The verifier reconciles nine registry/core/storage relations. The inference
fields are feature count $F$, total tree count ($T$ or $KR$), depth $d$,
scalar base ($a_0$ for v1, zero for the v2 reserved-field convention), and
scale $Q$. The storage fields are table pointer $P$, total bytes $\ell$,
number of chunks $N$, and stride $c$. For v2 the actual base vector comes
from the bytes after the header. This is a sufficient verification predicate;
not every listed field is consumed by every endpoint's arithmetic.

A conforming evaluator decodes the specified little-endian signed and
unsigned values, follows Equation~\ref{eq:traversal}, and computes the finite
integer sums without overflow or precision loss. A conforming local evaluator
first validates the canonical core. A conforming EVM evaluator obtains its
primitives from the manifest and its runtime metadata from the reconciled
record. Endpoint authorization, enabled-state, and payment conditions must be
satisfied so that execution reaches this evaluator with the unchanged packed
input.

\begin{theorem}[Conditional local/EVM execution agreement]
\label{thm:equivalence}
Let $B$ be canonical, let $M$ be well formed for $B$ at $(\sigma,r)$, and let
the nine registry/core/storage relations hold. Given the same $4F$ packed
bytes for $q\in\Ithirtytwo^F$, conforming local and EVM evaluators return the
same scalar or vector in Equation~\ref{eq:prediction}. Their class endpoints
return the same lowest-index maximizer and maximal logit.
\end{theorem}
\begin{proof}
Write a primitive offset as $o=jc+s$, with $0\le s<c$. If its length $n$ ends
within the chunk, the reader returns exactly $B[o:o+n]$. Otherwise it joins
$c-s$ bytes from chunk $j$ and the remaining $n-(c-s)$ bytes from chunk $j+1$.
Since $n\le4\le c$, a third chunk is unnecessary. Exact manifest framing and
lengths therefore supply identical primitive bytes to both evaluators.
Registry agreement supplies the same metadata where the contract obtains it
from the record rather than the core header.

At each root both evaluators start at index zero. If their indices agree
before a decision, they read the same feature and threshold and apply the
same strict comparison, so their next indices agree. Induction over $d$
levels selects the same leaf. Indeed, after $j$ decisions,
$2^j-1\le i_j\le2^{j+1}-2$, so subtracting $I(d)$ at level $d$ yields a valid
leaf index. Repeating over the finite ensemble gives identical summands.
Exact accumulation yields identical outputs, and the shared tie rule yields
identical class results.
\end{proof}

The theorem establishes an implication between specified objects and
algorithms. It does not prove that arbitrary Solidity bytecode conforms, that
an RPC returned the authentic state, or that a named record contains its
claimed bytes. The tests below address enumerated implementation paths.
The supplement gives the full geometry, traversal, reconstruction, and
quantization proofs without treating their elementary components as separate
algorithmic inventions.

\subsection{Content binding and a counterexample}

Execution agreement for reconstructed bytes does not require their registered
name to be their digest. If the name is claimed as a content identifier, the
verifier must additionally establish
\begin{equation}
 \mathrm{modelId}=\operatorname{keccak256}(B).
 \label{eq:binding}
\end{equation}
The deployed registry accepts the identifier and storage tuple as separate
arguments and does not reconstruct and hash the core during registration.
An admissible nonzero unused identifier different from the core digest is
therefore a counterexample to unconditional binding.

The isolated-EVM test instantiates this case against the unmodified contracts.
A valid 104-byte v1 object was registered under an identifier different from
its computed Keccak-256 digest. The runtime followed that record and returned
$1043$, the referenced object's prediction. The production JavaScript loader
reconstructed the same bytes, detected the mismatch, and rejected the record.
The complete identifiers and executable fixture are included in the
supplement and integration record. This is a client-enforced verification
condition; it does not introduce an invariant into already deployed bytecode.

This distinction generalizes beyond a particular tree format. A registry
lookup, a content digest, and execution under metadata are different links in
a computation's provenance. A positive observation for all records at one
block cannot eliminate a counterexample allowed by the registration
interface. Verification must check each link for the record being used.

\subsection{Bounds and approximation}

Scalar inference performs exactly $Td$ decisions and $T$ leaf reads; vector
inference performs $KRd$ decisions and $KR$ leaf reads. An abstract streaming
evaluator needs constant scalar traversal state or $O(K)$ output state,
excluding its inputs and read buffers. Finite work does not imply execution
within a particular gas limit. Moreover, the current Solidity boundary reader
allocates fresh buffers for crossing primitives; EVM memory can grow during
a call even though each primitive is small.

Under registry/core agreement, the count bound makes scalar sums and every
partial sum no larger
in magnitude than $(65535+1)2^{31}=2^{47}$. Since $KR\le65535$, each vector
component is bounded at least as tightly. These are below JavaScript's
$2^{53}$ exact-integer boundary and far below int256 limits. They apply to the
documented deployment profile, not arbitrarily large external files.

For a finite real value $z$ and positive integer $Q$, the abstract quantizer is
\begin{equation}
 {\cal Q}_Q(z)=\operatorname{clamp}_{\Ithirtytwo}
 \left(\left\lfloor Qz+\tfrac12\right\rfloor\right).
 \label{eq:quantizer}
\end{equation}
It rounds half ties toward positive infinity. Without saturation, the
single-value error is at most $1/(2Q)$. If quantizing a base and the selected
leaves leaves every path unchanged, scalar output error is at most
$(T+1)/(2Q)$; replace $T$ by $R$ for one vector component. A split margin
$|x-\tau|>1/Q$ suffices to preserve that split under this abstract quantizer.
These elementary bounds concern quantization of a fixed forest, not the error
between separately trained models.

Production preprocessing first computes a binary64 product
$\operatorname{fl}(Qz)=Qz+\delta_Q(z)$ and then rounds. Its unsaturated
single-value bound is $(1/2+|\delta_Q(z)|)/Q$. The abstract half-unit bound
therefore requires exact agreement with the abstract quantizer or this
product-rounding correction. Saturation and a changed path invalidate the
stated stable-path bound; a small feature perturbation near a threshold can
select a substantially different leaf.

\section{Implementation and experimental methods}
\label{sec:methods}

\subsection{Training and conformance profile}

The three trainers support regression, binary, multiclass, and multilabel
objectives with histogram-based fixed-depth trees. They use seeded
\codeword{xorshift32}, linear or quantile bins, explicit feature and output
ordering, and shared serialization. When a split is unavailable, a complete
continuation duplicates the terminal value and uses feature zero with the
maximal int32 threshold. It preserves the terminal result and fixed-depth
geometry for every packed input.

Training uses floating-point gradients, split scores, metrics, exponentials,
and logarithms. The publication profile fixes a finite float32 input matrix,
row and feature order, target representation, hyperparameters, seed,
random-consumption order, and numeric operation sequence. The C++ build uses
\codeword{-ffp-contract=off -fno-fast-math}. These restrictions allow
meaningful byte-level comparisons without asserting that every platform or
mathematical library will emit the same model. Comparisons exclude optional
JSON metadata containing timestamps.

Three observed discrepancies motivate the profile. First, a standalone linear-binning
witness selects bin four using division by the range and bin three using
multiplication by a rounded cached reciprocal, although the expressions are
algebraically equivalent. Second, a generated
4,000-row, 12-feature, four-class fixture exposed a one-byte core difference
when the softmax denominator accumulated binary32-rounded rather than
binary64 exponentials. Third, evaluating $\lfloor z+1/2\rfloor$ with a
binary64 addition rounded a value immediately below a half boundary upward
before the floor. The Python and C++ rounding implementations now compare the
fractional part with $1/2$, matching JavaScript's finite rounding convention.
The operation order, softmax precision, half-boundary, and explicit seed-zero
cases are retained as regression controls. These corrections close the
observed failures and remain empirically testable.

\subsection{Evaluation design}

The evaluation separates four questions: whether implemented parsers and
evaluators agree at format and arithmetic boundaries; whether the three
trainers emit identical cores for declared profiles; how identical ensembles
compare under two EVM storage implementations; and whether a public deployment
can be reconstructed and its reported outputs replayed from an archived state
observation. The synthetic fixtures target execution behavior and contain no
claim about a predictive application.

The parity dataset has 240 rows and five mixed-scale finite features, with
deterministically generated regression, binary, three-class, and three-label
targets. Baseline profiles use 18 depth-3 trees, learning rate 0.075, minimum
leaf size four, 16 bins, $Q=100000$, a 70/20/10 split, and seed 20260724.
The tests compare entire canonical core byte strings, rather than only
approximate outputs. The checking parser in
\path{tests/publication/model_format.py} and the bigint inference oracle are
separate implementations from the production decoder.

The local EVM experiments compile the production contracts with Solidity
0.8.20, optimizer 200, \codeword{viaIR=true}, and
\codeword{evmVersion=istanbul}.\cite{solidity0820} They execute in Ganache
7.9.2 under its recorded hardfork configuration.\cite{ganache792} The
production runtime compiles to 17,794 runtime bytes, below EIP-170. The
compiler profile is deliberately pinned because a permissive source pragma
is not evidence that all compiler versions build the same program.

The recorded CPU study uses Python 3.12.13, NumPy 2.3.5,\cite{harris2020numpy}
Node 24.19.0, and g++ 13.3.0 on Linux x86-64 with an AMD EPYC 9V74 host and
nine logical CPUs visible. Its processes are single-threaded, and the host
was not isolated or CPU-pinned. Detailed environment and source digests are
stored with each result; new evidence retains its own collection provenance.

\section{Results}
\label{sec:results}

\subsection{Training and malformed-artifact conformance}

All 25 distinct three-engine training profiles produced identical core bytes.
The record also contains five auxiliary controls, for 30 profile executions,
and one separately counted arithmetic operation-order witness. The 18 test
methods exercise all four objectives, both binning methods, weighting,
stratification, refitting, learning-rate schedules, repeated execution, and
boundary inference. This is a finite matrix, not a branch-coverage percentage
or a statistical estimate of cross-platform reliability.

The enabled early-stopping and plateau configurations in the frozen matrix
do not actually trigger early termination or a plateau reduction. They verify
option propagation and continued-training parity; they do not cover those
positive control-flow branches. Making that distinction avoids counting a
configuration as evidence for behavior it did not execute.

The malformed-artifact tests exercise unknown versions, nonzero reserved
fields, inconsistent dimensions, nonpositive scale, out-of-range features,
truncation, trailing bytes, metadata framing, and non-finite inputs. Eight
additional executable property tests exercise traversal intervals,
root-to-leaf paths, one-/two-/four-byte reads across nine chunk sizes,
a three-chunk counterexample when $c<4$, pointer capacity, rational
quantization bounds, accumulator bounds, and checked-in artifact equality.
All eight passed. The resulting evidence supports the named failure modes
rather than exhaustive validation of every malformed byte string.

\subsection{Local EVM boundary tests}

The integration test uses an isolated chain. It deploys the production
store, registry, model NFT, and inference runtime, then
publishes a 104-byte scalar fixture and a 156-byte three-output fixture using
17-byte chunks, deliberately forcing primitive reads across boundaries.
Readback reconstructs both cores exactly and verifies their content digests.
Six packed vectors exercise threshold equality, neighboring branches, signed
int32 extrema, and equal-logit class selection. The test also checks rejection
of invalid input length.

The separate bigint evaluator matches 18 scalar/vector/class view results and
two local transaction results, with zero mismatches. Equal logits select class
zero as specified. The registration counterexample in
Section~\ref{sec:assurance} is exercised in the same run. These transactions
are local test-chain operations; no experiment changes the public deployment.

A scalar scaling experiment publishes six deterministic two-feature shapes
through 20,000-byte chunks and evaluates 12 seeded vectors per shape. All
72 comparisons agree with the bigint oracle. Core sizes range from 1,784 to
38,024 bytes. Recorded mean view-gas estimates range from 286,247 for 20
depth-3 trees to 3,121,210 for 200 depth-4 trees. Fifty depth-6 trees occupy
more bytes than 200 depth-4 trees but require fewer decisions and a lower
mean estimate, illustrating why byte size alone is insufficient to describe
inference cost. The supplement retains the full scaling table; the controlled
comparison below evaluates the storage decision directly.

\subsection{Controlled storage comparison}
\label{sec:storagecomparison}

The comparison stores the same six scalar cores both through the production
code-as-data path and through a benchmark-only Solidity
\codeword{mapping(bytes32 => bytes)}. The latter packs the canonical bytes
into storage words and reads each primitive through an indexed
\codeword{SLOAD}. Both paths use the same registered model, two matching
registry lookups, free-view access conditions, feature vector, traversal rule,
and exact accumulation. They share the compiler and Ganache/Shanghai
configuration. The code path uses 20,000-byte chunks.

The experiment measures two costs separately. Per-model publication is the
sum of mined local receipts for chunk and table writes, or the corresponding
packed-storage write. Common registry registration is reported separately in
the result record. Inference is measured with one \codeword{eth\_estimateGas}
request per path and input, resetting transaction access warmth for each
request. Shared contract deployment is also recorded separately: the
benchmark-only comparator has a smaller feature surface, so its deployment
cost cannot support a like-for-like system-cost claim.

All 72 model/input combinations produce the same result in the bigint oracle,
production runtime, and packed-storage evaluator: 144 reference/EVM
comparisons with zero mismatches. Readback also verifies identical stored core
bytes. Every production inference estimate reproduces the corresponding
observation in the earlier scalar scaling record.

\begin{table}[htbp]
\centering
\small
\caption{Controlled local-EVM storage comparison. Publication values are sums
of mined write receipts; inference values are arithmetic means of 12 gas
estimates. Code: production code-as-data implementation. Storage: packed
Solidity-storage comparator. Common registry and shared deployment costs are
excluded. Rounded values are shown; raw observations are archived.}
\label{tab:storagecomparison}
\setlength{\tabcolsep}{4pt}
\begin{tabular}{@{}rrrrrr@{}}
\toprule
 & & \multicolumn{2}{c}{Publication ($10^6$ gas)} & \multicolumn{2}{c}{Inference ($10^3$ gas)}\\
$T,d$ & Core bytes & Code & Storage & Code & Storage\\
\midrule
20, 3 & 1,784 & 0.500 & 1.310 & 286.2 & 267.1\\
50, 3 & 4,424 & 1.062 & 3.184 & 644.5 & 597.7\\
50, 4 & 9,224 & 2.078 & 6.564 & 810.8 & 783.9\\
100, 4 & 18,424 & 4.003 & 13.012 & 1576.9 & 1518.8\\
200, 4 & 36,824 & 7.951 & 25.962 & 3121.2 & 2992.5\\
50, 6 & 38,024 & 8.204 & 26.817 & 1146.4 & 1149.5\\
\bottomrule
\end{tabular}
\end{table}

Across these shapes, code-as-data publication consumes 61.8--69.4\% less gas
than the packed-storage write, equivalently a storage-to-code gas ratio of
2.62--3.27. The materialized byte footprints are nearly equal: code accounting
includes chunk prefixes and the pointer table, while storage accounting
includes the length slot and rounded-up data slots. Both exclude common
registry metadata, executable code, account records, and trie overhead.
The observed write saving is therefore a gas-schedule and implementation
result, rather than a compression result.

Inference shows a different tradeoff. Packed storage uses approximately
3.3--7.3\% less gas for five shapes and 0.27\% more for the depth-6 shape.
The production reader repeatedly resolves table entries and checks framing;
the comparator reads packed words directly. Thus this is a comparison of
implemented storage backends, not an isolated comparison of opcode costs.
The results support lower model-publication gas for the measured code-as-data
construction, but no universal inference-gas advantage. They cover one
synthetic scalar family, one client, one compiler profile, and one chunk
size; vector models and alternative optimized layouts remain unmeasured.


\subsection{Replayable deployment observation}
\label{sec:deployment}

A fresh read-only observation selected GenesisL1 chain ID 29, block
13,602,838, timestamp 6 September 2026 at 17:20:41 UTC, and block hash
\begin{center}
\small\hashpair{0xfd3da1020c37ee3c1fe7cd0a6060dbc5e}{c3ec5fb90c0b812256dc33e467dace3}.
\end{center}
The witness and archive collector queried \url{https://rpc.genesisl1.org}.
Every state-dependent getter, code read, and inference call used the selected
numeric block tag. Each collection stage re-read the block hash after its
calls and required it to remain unchanged. The collector additionally
required the explicitly selected number/hash pair before accepting state.
No public transaction was submitted.

The observer enumerated all 12 active token records, recovered their address
tables and 1,306 data chunks, and reconstructed 31,185,324 core bytes. Every
core passed strict format checks, all nine registry/core/storage relations,
and the registered Keccak-256 commitment. The largest object contains
9,523,096 bytes in 397 chunks. Table~\ref{tab:deployment} records the
per-model geometry. All models use a nominal 24,000-byte chunk stride and
were enabled for unrestricted view inference at the selected state.

\begin{table}[htbp]
\centering
\small
\caption{Public model geometry at block 13,602,838. Trees denotes total
serialized trees ($T$ or $KR$). Token identifiers locate the observed records;
registry titles and application claims are not used as scientific labels.
Every row passed content, canonical-format, manifest, and nine-relation
verification.}
\label{tab:deployment}
\begin{tabular}{@{}rrrrrrr@{}}
\toprule
Token & Version & $F$ & $d$ & Trees & Core bytes & Chunks\\
\midrule
1 & 2 & 4 & 3 & 1,800 & 158,436 & 7\\
2 & 1 & 7 & 5 & 169 & 63,568 & 3\\
3 & 1 & 13 & 6 & 196 & 148,984 & 7\\
4 & 1 & 4 & 3 & 275 & 24,224 & 2\\
5 & 1 & 53 & 5 & 20 & 7,544 & 1\\
6 & 1 & 26 & 8 & 1,496 & 4,583,768 & 191\\
7 & 1 & 27 & 8 & 2,500 & 7,660,024 & 320\\
8 & 1 & 33 & 8 & 1,008 & 3,088,536 & 129\\
9 & 1 & 33 & 9 & 1,552 & 9,523,096 & 397\\
10 & 1 & 33 & 7 & 1,789 & 2,733,616 & 114\\
11 & 1 & 33 & 8 & 1,042 & 3,192,712 & 134\\
12 & 1 & 47 & 3 & 9 & 816 & 1\\
\bottomrule
\end{tabular}
\end{table}

For each model, nine deduplicated packed vectors comprise the all-zero vector,
both signed-int32 extrema, and complete $\theta-1$, $\theta$, $\theta+1$
neighborhoods for two distinct non-sentinel root thresholds. Other feature
coordinates remain fixed for each threshold probe. Roots are necessarily
visited, so the neighborhoods exercise their named branch decisions; they
do not cover every internal path. All 108 local/provider-returned scalar or
vector outputs matched exactly.

The compressed evidence archive contains the model cores, exact storage
runtime objects and address tables, five deployed contract runtimes, the
block records, raw RPC request/response transcript, packed vectors, expected
outputs, collection source, and file digests. A separate Python verifier uses
only the standard library and does not import the trainers, production
decoder, or JavaScript witness. It checks the file digests, parses the raw
observations, reconstructs each object, recomputes Keccak-256, validates the
registry relations, and evaluates every archived input with integer
arithmetic. Offline replay verified all 12 models, all 31,185,324 core bytes,
and all 108 output comparisons without an RPC connection.

The archive is distributed as
\path{benchmarks/results/live_chain_archive_13602838.tar.gz}, with a small
machine-readable result and checksum alongside it. Extraction and verifier
commands are provided in \path{REPRODUCIBILITY.md}. Its raw observations
support inspection and replay of what the provider returned. Neither the
transcript nor the offline verifier authenticates consensus finality,
validates account-proof paths, or establishes source-to-live-bytecode
identity. Contract runtime bytes and their digests are retained as observed
objects, not as a verified compilation claim.

This is a new September observation. Earlier July summaries at block
13,342,043 remain available as historical records, but their missing raw
responses are not reconstructed retrospectively or used as the principal
deployment evidence. The present result is also bounded to one provider and
one state. It does not imply that later registrations satisfy the predicate
or that the 12 named applications have been predictively validated.


\subsection{Reference trainer timing}

The existing end-to-end timing study uses 3,000 rows, 12 deterministic
features, and 60 depth-4 trees for regression and binary classification. Each
engine receives one warm-up and 30 timed repetitions in rotating order.
Timing includes startup, parsing, training, serialization, and output writing.
Python and C++ read CSV, while Node reads an equivalent JSON matrix; the
measurements therefore describe complete workflows with different parsers.

Regression medians are 0.3046 s for Python, 0.0365 s for C++, and 0.1384 s for
Node; binary medians are 0.3149 s, 0.0437 s, and 0.1610 s respectively. Both
workloads produce identical 11,064-byte cores across all three engines.
The supplement reports interquartile ranges, and the JSON record retains all
measurements. These observations establish within-host reference performance;
they do not establish a general language ranking or compare GL1F with
optimized external training libraries.

\section{Discussion and limitations}
\label{sec:discussion}

GL1F makes a limited connection between public storage and reproducible
computation explicit. Canonical decoding establishes the represented forest;
ordered reconstruction establishes which bytes are stored; registry checks
establish which metadata the runtime uses; and shared integer semantics
establish the returned function. A digest binds the requested identifier to
those bytes only when it is actually recomputed. The executable registry
counterexample is useful evidence of this separation: correct execution of
some referenced bytes can coexist with an incorrect content identifier.

Complete fixed-depth trees trade storage for regular control flow and direct
offsets. Forced continuations can duplicate terminal values, and the current
internal-node layout contains two reserved bytes per node. These choices
simplify interpretation but do not minimize serialization size. Alternative
sparse encodings, compressed nodes, or batched reads could change the cost
balance. The storage experiment evaluates a specific packed-storage baseline
under a fixed EVM configuration; it is not a proof of optimality against all
possible implementations.

The formal result is conditional and handwritten. It is supported by
source inspection and separate executable oracles, but the study does not
prove compiler correctness or mechanically verify deployed bytecode. Positive
tests cannot establish universal parser safety or training parity. In
particular, differing transcendental libraries, operation reassociation, or
unexercised stopping branches can change a trained core. The finite fixture
matrix and retained counterexamples make those boundaries concrete.

The fresh deployment archive improves reproducibility by preserving the
inputs to verification. Offline replay can detect a corrupt archive, an
inconsistent manifest, a content mismatch, or disagreement between recorded
and locally computed outputs. It cannot establish that the provider supplied
an authentic header or account state. Such a claim would additionally need
an authenticated chain view and verified state proofs. A block-number/hash
recheck detects a changed response during collection but does not remove the
provider trust assumption.

The case study is developer-conducted and concerns one deployment. Separately
implemented verification code reduces shared implementation dependence but
is not independent authorship or external replication. Creator addresses and
model titles do not establish independent adoption, dataset provenance, or
predictive usefulness. Application-level evaluation would require frozen data,
leakage-resistant splits, suitable baselines, uncertainty estimates, and
external validation; the synthetic conformance experiments do not supply
those results.

The public-byte design also limits what ownership and payment mechanisms can
promise. Registry settings, fee destinations, authorization keys, and token
ownership can evolve separately from the core. Their correctness must be
checked for the service being used, rather than inferred from the execution
theorem. A scientific result should identify both the fixed artifact and the
state at which its service behavior was observed. None of the reported
experiments requires a successor contract or migration of existing models.

\section{Conclusion}

GL1F provides a reproducible method for relating a registered public tree
ensemble to the integer function it executes. Its central requirement is a
composition of canonical parsing, exact ordered reconstruction, consistent
runtime metadata, identical packed inputs, and exact arithmetic. The
specification proof explains why these conditions suffice; the registry
counterexample explains why an identifier lookup alone does not. Cross-language
conformance, local EVM boundary tests, controlled storage measurements, and a
replayable deployment archive provide evidence at distinct levels. The result
is a bounded execution-assurance study whose checks can be repeated and whose
remaining assumptions are explicit.

\ledgernotes
\section*{Data and software availability}

The MIT-licensed reference implementation, canonical format specification,
executable fixtures, benchmark methods, recorded results, and manuscript
sources are available in the GenesisL1 Forest repository.\cite{gl1frepository}
Third-party material retains its stated upstream terms. The supplement
contains detailed proofs, field layouts, and supplementary measurements.
\path{REPRODUCIBILITY.md} gives the pinned environment and commands. The
aggregate \codeword{make verify} command runs local software checks;
\codeword{make pdfs} builds the manuscript and supplement. The storage
comparison and deployment archive each include the provenance required to
replay their reported checks.

\section*{Competing interests}

Mikhail Fedorov is the founder of GenesisL1 and the developer of GL1F. The
paper evaluates software and a blockchain deployment associated with that
project. This relationship is a potential competing interest.

\section*{Use of artificial intelligence}

OpenAI ChatGPT/Codex assisted with manuscript drafting and revision, code
inspection and corrections, test and benchmark development, and preparation
of reproducibility materials. Generated content was checked through source
comparison, executable tests, and numerical and document review. These checks
are described with their scope and limitations in the paper and accompanying
materials. The author is responsible for the submitted manuscript and its
scientific claims.

\bibliographystyle{ledgerbib}
\bibliography{references}
\end{document}
